\documentclass[12pt]{article}
\usepackage[utf8]{inputenc}
\usepackage[T1]{fontenc}
\usepackage{hyperref}
\title{Yango Language Resources for K-12 Curriculum Development}
\author{Compiled by Claude Sonnet 5.5 (claude-sonnet-5-5) for Andr\'as Kornai}
\date{}
\begin{document}
\maketitle

\section{Introduction}
Yango (ISO 639-3 yng; Glottocode yang1289; also called Gbendere) is a Ubangian language of the Ngbaka group, spoken in the Democratic Republic of the Congo. The search placed speakers in Equateur Province, in the Kungu and Libenge districts near the Esobe River. Estimates of about 3,000 speakers come from Verbix and ASJP. Its relation to Monzombo and Kpala is unclear, since they are listed as separate varieties. Note that the reference description in the task calls the family Atlantic-Congo, while the sources found treat it as Ubangian. Very little educational material turned up.

\section{Primary Educational Resources}
No primers, textbooks or school programmes were found.

\section{Linguistic and Reference Materials}
\begin{itemize}
\item Wikipedia, Ngbaka languages (classification context): \url{https://en.wikipedia.org/wiki/Ngbaka_languages}
\item Wikipedia, Monzombo language (closely related variety): \url{https://en.wikipedia.org/wiki/Monzombo_language}
\item ASJP word-list data for Yango: \url{https://asjp.clld.org/languages/YANGO}
\item Endangered Languages Project entry, with a bibliography page: \url{https://endangeredlanguages.com/lang/4600/bibliography}
\end{itemize}

\section{Religious and Community Resources}
None found.

\section{Digital and Audio Resources}
\begin{itemize}
\item Verbix endangered languages page for Yango: \url{https://docs.verbix.com/EndangeredLanguages/Yango}
\end{itemize}
No audio, video or software was found.

\section{Curriculum Development Context}
The language is not known to be taught in schools. Documentation appears to be limited to word lists and survey-level entries. A curriculum would first need community-led documentation. It could also draw on related Ngbaka-group materials or on the regional language of schooling, but this is a suggestion and not something the search showed.
\end{document}